\documentclass[11pt,letterpaper]{article}
% isms-27000-style.tex
% Suite preamble for the ISO/IEC 27000/27001/27002 ISMS engineering suite.
%
% RECONSTRUCTION PROVENANCE: the ten uploaded suite documents load a shared
% preamble with `\input{style.tex}`, but no `style.tex` was present in the
% upload batches or anywhere in this repository's history. This file is a
% documented, suite-specific reconstruction inferred from the macros, colors
% and environments the documents actually use (colors navy/teal/graytext/pale;
% tag commands \SourceTag, \EngineeringTag, \SourceNote{..}; breakable
% tcolorbox callouts; color-striped longtables; TikZ flow diagrams). It is
% NOT the recovered original style; the original visual design is unknown.
%
% pdfLaTeX only.
\usepackage[T1]{fontenc}
\usepackage[utf8]{inputenc}
\usepackage{lmodern}
\usepackage{microtype}
\usepackage[letterpaper,margin=0.9in,headheight=15pt]{geometry}
\usepackage[table]{xcolor}
\definecolor{navy}{HTML}{14355A}
\definecolor{teal}{HTML}{0A7580}
\definecolor{graytext}{HTML}{55616E}
\definecolor{pale}{HTML}{EFF5F7}
\usepackage{amsmath,amssymb}
\usepackage{array,booktabs,longtable,tabularx,multirow}
\usepackage{enumitem}
\usepackage[most]{tcolorbox}
\usepackage{tikz}
\usetikzlibrary{arrows.meta,positioning,shapes.geometric,fit,calc}
\usepackage{needspace}
\usepackage{titlesec}
\usepackage{fancyhdr}
\usepackage{lastpage}
\usepackage{hyperref}
\hypersetup{colorlinks=true,linkcolor=navy,urlcolor=teal,citecolor=navy,pdfauthor={ISMS Engineering Documentation Suite},pdfsubject={Original supporting guidance for ISO/IEC 27000, 27001 and 27002}}
\setlength{\parindent}{0pt}
\setlength{\parskip}{5pt plus 1pt}
\setlength{\emergencystretch}{2em}
\renewcommand{\arraystretch}{1.18}
\setlist[itemize]{itemsep=2pt,topsep=3pt,leftmargin=1.5em}
\setlist[enumerate]{itemsep=2pt,topsep=3pt,leftmargin=1.7em}
\titleformat{\section}{\Large\bfseries\color{navy}}{\thesection}{0.7em}{}[\vspace{-3pt}\color{teal}\titlerule]
\titleformat{\subsection}{\large\bfseries\color{navy}}{\thesubsection}{0.6em}{}
\titleformat{\subsubsection}{\normalsize\bfseries\color{teal}}{\thesubsubsection}{0.55em}{}
\pagestyle{fancy}
\fancyhf{}
\fancyhead[L]{\sffamily\footnotesize ISO/IEC 27000 series | ISMS Engineering}
\fancyhead[R]{\sffamily\footnotesize\color{graytext} Independent companion}
\fancyfoot[L]{\sffamily\scriptsize Original explanatory material; consult the licensed standards}
\fancyfoot[R]{\sffamily\scriptsize\thepage\ / \pageref{LastPage}}
\renewcommand{\headrulewidth}{0.3pt}
% Inline provenance tags used throughout the suite: \SourceTag marks
% statements grounded in the standards; \EngineeringTag marks original
% engineering interpretation. Both are zero-argument label commands.
\newcommand{\SourceTag}{\textsf{\footnotesize\bfseries\color{navy}[Standard]}~}
\newcommand{\EngineeringTag}{\textsf{\footnotesize\bfseries\color{teal}[Engineering]}~}
\newcommand{\SourceNote}[1]{\par\smallskip{\footnotesize\color{graytext}\textit{Trace: #1.}}\par}

\begin{document}
\pagenumbering{roman}
\begin{titlepage}
\vspace*{1.2in}
{\color{teal}\rule{\textwidth}{4pt}}\par\vspace{20pt}
{\Huge\bfseries\color{navy} Governance and Risk Management}\par\vspace{18pt}
{\Large\color{graytext} Decision rights, scenario-based risk treatment, SoA governance and executive review}\par\vspace{1.3in}
{\large ISO/IEC 27000:2018 $\vert$ ISO/IEC 27001:2022 $\vert$ ISO/IEC 27002:2022}\par
\vfill
\textbf{Independent educational and engineering implementation reference}\par
Version 1.0 $\cdot$ October 2026\par
{\small Original explanatory material. Source standards are copyrighted; consult licensed originals for authoritative requirements.}\par
\end{titlepage}
\tableofcontents\clearpage\pagenumbering{arabic}
\section{Governance design and decision rights}
\SourceTag ISO/IEC 27001 clauses 4--6 establish context, leadership, risk planning and objectives; ISO/IEC 27000:2018 supplies terminology for governance, risk and effectiveness. \EngineeringTag Use a named executive sponsor, accountable risk owners and a separate ISMS process owner.
\small\rowcolors{2}{pale}{white}\begin{longtable}{p{.34\linewidth}p{.27\linewidth}p{.33\linewidth}}
\rowcolor{navy} \textcolor{white}{\textbf{Decision}} & \textcolor{white}{\textbf{Accountable role}} & \textcolor{white}{\textbf{Minimum record}}\\
\endfirsthead
\rowcolor{navy} \textcolor{white}{\textbf{Decision}} & \textcolor{white}{\textbf{Accountable role}} & \textcolor{white}{\textbf{Minimum record}}\\
\endhead
ISMS scope and risk criteria & Top management / delegated governance & Scope approval and criteria versions\\
Risk treatment and residual exposure & Named risk owner & Dated acceptance with rationale\\
Control ownership and implementation & Control owner & Design, operation and evidence schedule\\
Control independence assessment & Internal audit lead & Auditor conflict declaration\\
Management system improvement & Top management & Review actions with due dates\\
\end{longtable}\normalsize
\section{Governance operating rhythm}
\small\rowcolors{2}{pale}{white}\begin{longtable}{p{.23\linewidth}p{.29\linewidth}p{.42\linewidth}}
\rowcolor{navy} \textcolor{white}{\textbf{Cadence (example)}} & \textcolor{white}{\textbf{Meeting or activity}} & \textcolor{white}{\textbf{Required input and output}}\\
\endfirsthead
\rowcolor{navy} \textcolor{white}{\textbf{Cadence (example)}} & \textcolor{white}{\textbf{Meeting or activity}} & \textcolor{white}{\textbf{Required input and output}}\\
\endhead
Weekly & Treatment backlog review & High-exposure risks, overdue actions, resource blockers\\
Monthly & Control effectiveness review & Performance trends, exceptions, incidents, change events\\
Quarterly & Enterprise risk and supplier review & Risk re-ratings, key supplier exposure, SoA modifications\\
Planned interval & Internal audit and management review & Independent findings, strategy and resource decisions\\
Trigger-based & Material business/technology change & Reassessment and approved ISMS change\\
\end{longtable}\normalsize

\section{Risk identification: write scenarios}
\begin{enumerate}
\item Identify \textbf{in-scope information and business process}. Link supporting cloud services, people, sites and suppliers.
\item Identify \textbf{causal conditions}: exploitable defect, weak access rule, unprotected key, loss of utility or human error.
\item State \textbf{threat event and effect}, e.g. stolen CI token permits source alteration leading to unauthorized production deployment.
\item Assign a \textbf{business risk owner}, CIA impact values and evidence supporting likelihood.
\item Compare to published acceptance criteria; consider treatment and transfer only as supported by real agreements.
\end{enumerate}

\section{Illustrative ordinal scoring model}
\EngineeringTag Score impact $I=\max(C,I_n,A)$ from separate confidentiality, integrity and availability ratings in $\{1,\ldots,5\}$; multiply by likelihood $L\in\{1,\ldots,5\}$ to obtain $R=L\cdot I$. Use this as \emph{ordinal prioritization}, not as a monetary expectation or a statistically calibrated probability. Document inter-rater calibration and material context exceptions.

\[R_{\mathrm{inherent}}=L_{\mathrm{before}}\,\max(C,I_n,A),\qquad R_{\mathrm{residual}}=L_{\mathrm{after}}\,I_{\mathrm{after}}.\]
\small\rowcolors{2}{pale}{white}\begin{longtable}{p{.09\linewidth}p{.37\linewidth}p{.11\linewidth}p{.37\linewidth}}
\rowcolor{navy} \textcolor{white}{\textbf{Impact}} & \textcolor{white}{\textbf{Working interpretation}} & \textcolor{white}{\textbf{Likelihood}} & \textcolor{white}{\textbf{Working interpretation}}\\
\endfirsthead
\rowcolor{navy} \textcolor{white}{\textbf{Impact}} & \textcolor{white}{\textbf{Working interpretation}} & \textcolor{white}{\textbf{Likelihood}} & \textcolor{white}{\textbf{Working interpretation}}\\
\endhead
1 & Negligible operational consequence & 1 & Rare under present conditions\\
2 & Limited team-level loss & 2 & Unlikely but credible\\
3 & Material service, contractual or data effect & 3 & Plausible\\
4 & Major sustained loss or serious regulatory exposure & 4 & Likely without change\\
5 & Severe enterprise-wide consequence & 5 & Expected or repeatedly seen\\
\end{longtable}\normalsize
\small\rowcolors{2}{pale}{white}\begin{longtable}{p{.15\linewidth}p{.49\linewidth}p{.30\linewidth}}
\rowcolor{navy} \textcolor{white}{\textbf{Score}} & \textcolor{white}{\textbf{Suggested handling}} & \textcolor{white}{\textbf{Decision authority}}\\
\endfirsthead
\rowcolor{navy} \textcolor{white}{\textbf{Score}} & \textcolor{white}{\textbf{Suggested handling}} & \textcolor{white}{\textbf{Decision authority}}\\
\endhead
1--4 & Track within normal control program & Risk owner\\
5--9 & Time-bound mitigation or documented acceptance & Risk owner with security consultation\\
10--15 & Prioritize treatment with milestones & Senior business owner\\
16--25 & Escalate rapidly and restrict exposure when feasible & Executive risk governance\\
\end{longtable}\normalsize
\section{Risk treatment and Statement of Applicability link}
\begin{enumerate}
\item Choose a documented response (e.g., mitigate, avoid, share or accept) appropriate to risk criteria.
\item Determine necessary controls, including organization-defined controls outside Annex A when necessary.
\item Compare determined controls with the 93 Annex A references; identify omissions or explicit exclusions.
\item Create and maintain an SoA recording necessary controls, inclusion reasoning, implementation state and exclusion justifications.
\item Obtain acceptance from the relevant risk owners for plans and residual exposure; capture due dates and proof of implementation.
\end{enumerate}

\section{Sample risk record}
\small\rowcolors{2}{pale}{white}\begin{longtable}{p{.26\linewidth}p{.68\linewidth}}
\rowcolor{navy} \textcolor{white}{\textbf{Attribute}} & \textcolor{white}{\textbf{Example}}\\
\endfirsthead
\rowcolor{navy} \textcolor{white}{\textbf{Attribute}} & \textcolor{white}{\textbf{Example}}\\
\endhead
Risk ID & R-APP-014\\
Asset / objective & Customer service API / authorized record changes\\
Scenario & Compromised deployment token enables unauthorized code delivery\\
Impact / likelihood & Integrity 5; likelihood 3; inherent level 15\\
Planned treatment & Constrain token privileges; protected builds; rotating credentials; deployment attestations\\
Control candidates & 5.16, 5.18, 8.2, 8.4, 8.24, 8.25, 8.32\\
Residual state & Reassess only after control validation, record evidence and approval\\
Owner and approver & Application product owner; appropriate delegated risk approver\\
\end{longtable}\normalsize

\section{Governance metrics and warnings}
\begin{itemize}
\item Risk treatment closure rate: actions verified on or before approved due date / all due actions; exclude unverified ticket closures.
\item Evidence freshness: applicable control tests with valid sampling period / controls planned for test; track N/A separately.
\item Residual acceptance quality: approved risk decisions containing owner, rationale, expiration and evidence / accepted risks.
\item Change sensitivity: material ISMS scope changes with documented risk reassessment / all such changes.
\item Avoid combining operational security scanner coverage with compliance assurance; one is an input to the other.
\end{itemize}

\section{Management review agenda}
\begin{enumerate}
\item Previous actions and ageing
\item Changes to context, stakeholders and obligations
\item Effectiveness indicators and security objectives
\item Audit findings, nonconformities and corrective outcomes
\item Risk assessment results and outstanding treatment
\item Resource conflicts, supplier dependencies and improvement decisions
\end{enumerate}

\clearpage\section*{Source and interpretation policy}
\addcontentsline{toc}{section}{Source and interpretation policy}
\textbf{Source editions:} ISO/IEC 27000:2018 (overview and vocabulary); ISO/IEC 27001:2022 (requirements and normative Annex A); ISO/IEC 27002:2022 corrected 2022-03 (control guidance). Specific clause identifiers are references to the supplied editions, not reproduced control text.\par
\textbf{Interpretation:} ``Source-grounded paraphrase'' describes a concept in the cited source. ``Engineering recommendation'' is a practical model, technology choice, owner, evidence example, frequency, template or scoring convention authored for this suite; it is not a claim that ISO prescribes that exact approach.\par
\textbf{Conformity:} Clauses 4--10 of ISO/IEC 27001 govern ISMS conformity. Applicability of controls is determined through risk treatment, with Annex A used as a completeness cross-check and exclusions justified in the Statement of Applicability (SoA). Source documents alone neither certify an ISMS nor prove control operating effectiveness.\par
\textbf{Source licensing:} Do not distribute or embed the licensed standard PDFs in repositories or distribute copied normative text. Review the editions, amendments, and certification scheme applicable to each assessment before relying on this work.\par
\end{document}
